\section{Conclusion}
\label{sec:conclusion}

\rev{We evaluated explainability-guided electrode reduction for subject-independent sEMG gesture recognition using cross-fitted attribution and prespecified primary decision rules.}

\rev{\textit{Supported}, for subject-independent, offline recognition with a two-ring monopolar forearm montage in able-bodied participants: removal of four of 16 electrodes within a prespecified 0.02 Macro-F1 tolerance, unchanged under Holm correction, but not of five; attribution-ranked subsets above the mean of random subsets of the same size; deletion-based support for the faithfulness of the ranking; and no measurable benefit from the stability penalty. In an exploratory, post-hoc replication on putEMG, $k^* = \PutKStarPtTwo$ of 24 electrodes met the same tolerance.}

\rev{\textit{Not established}: clinical performance; reduction of physical bipolar hardware (with derived bipolar pairs, no pair could be removed within the same tolerance); equivalence to, or superiority over, a data-free geometric selection; real-time use; muscle-level localisation; and generalisation beyond these cohorts and montages.}

\rev{Validation in target clinical populations (e.g.\ people with upper-limb loss) and in online, causal control settings is required.}
